\documentclass[10pt,a4paper]{amsart}
\usepackage[margin=19mm]{geometry}
\usepackage{amsmath,amssymb,mathtools}
\usepackage[scr=rsfs,cal=euler]{mathalfa}
\usepackage{xfrac}
\usepackage[unicode,hidelinks,bookmarks=false]{hyperref}
\newcommand{\ie}{i.e.\ }
\newcommand{\cf}{cf.\ }
\newcommand{\resp}{respectively\ }
\newcommand{\Subsection}{\subsection}
\providecommand{\nobreakdash}{\nobreak\hbox{-}}
\setlength{\emergencystretch}{2em}
\multlinegap0pt
\newtheorem{theorem}{Theorem}[section]
\newtheorem{thmy}{Theorem}
\renewcommand{\thethmy}{\Alph{thmy}} % "letter-numbered" theorems
\newenvironment{thmx}{\stepcounter{thm}\begin{thmy}}{\end{thmy}}
\newtheorem{lemma}[theorem]{Lemma}
\newtheorem{corollary}[theorem]{Corollary}
\newtheorem{conjecture}[theorem]{Conjecture}
\newtheorem{proposition}[theorem]{Proposition}
\newtheorem{example}[theorem]{Example}
\newtheorem{remark}[theorem]{Remark}
\newcommand{\calc}{{\cal C}}
\newcommand{\N}{\mathbb{N}}
\newcommand{\Z}{\mathbb{Z}}
\newcommand{\dd}{\displaystyle }
\newcommand{\nn}{\N}
\newcommand{\X}[2]{\stb{#1}{#2}{\mbox{\Huge$\times$}}}
\newcommand{\bld}[2]{{\buildrel{#1}\over{#2}}}
\newcommand{\st}[2]{{\mathrel{\mathop{#2}\limits_{#1}}{}\!}}
\newcommand{\stb}[3]{{\st{{#1}}{\bld{{#2}}{#3}}{}\!}}
\newcommand{\xmare}[2]{\stb{#1}{#2}{\mbox{\Huge$\times$}}}
\newcommand{\w}{\wedge }
\newcommand{\s}{{\sigma}}
\newcommand{\0}{{\leqno}}
\newcommand{\CD}{\mathcal{C}\mathcal{D}}
\def\barr{\begin{array}}
\def\earr{\end{array}}
\begin{document}
\title[Two new functions related to the sum of element orders of a ]{Two new functions related to the sum of element orders of a finite group}
\author{Marius Tărnăuceanu}
\date{}
\maketitle
\noindent\emph{Source and licence.} Two new functions related to the sum of element orders of a finite group. Version arXiv:2608.10035v1 (CC BY 4.0). This material is distributed under the Creative Commons Attribution 4.0 International licence, http://creativecommons.org/licenses/by/4.0/, as recorded in the arXiv OAI licence metadata for this exact version and shown on the abstract page https://arxiv.org/abs/2608.10035v1. Full rights notice: licenses/arxiv-2608.10035-CC-BY-4.0.txt.\par\smallskip
\noindent\textit{Editorial scope.} Counted unit: the original \texttt{lemma} environment \texttt{-} at source lines 163--172 together with its complete proof at lines 173--192; the delivered document keeps source lines 146--193 so that every referenced label and every citation resolves inside it. Native editable source is the paper's own concatenated TeX; the AMS wrapper is editorial formatting only. No publisher class, upstream script or unrelated artwork is redistributed.
\section{Proofs of the main results}

First of all, we study equal order pairs $(G,H)$ satisfying
\begin{equation}
o(x)=o(xH), \forall\, x\in G\setminus H,
\end{equation}and
\begin{equation}
o(x)=o(xH)|H|, \forall\, x\in G\setminus H,
\end{equation}respectively. These correspond to the cases when we have equality in the inequalities (1) and (2).

We observe that the condition (3) is equivalent to
\begin{equation}
\langle x\rangle\cap H=1, \forall\, x\in G\setminus H,\nonumber
\end{equation}i.e. to the fact that $H$ is an isolated subgroup. Clearly, $(C_p\times C_p,C_p\times 0)$ with $p$ prime is an equal order pair satisfying (3), while $(D_8,Z(D_8))$ and $(Q_8,Z(Q_8))$ are equal order pairs that does not satisfy (3).

An interesting example of such a pair is given by the following lemma. 


\begin{lemma}
Let $G={\rm ZM}(m,n,r)$ and $H=\langle a\rangle$. Denote by $d$ the multiplicative order of $r$ modulo $m$. Then the following conditions are equivalent: 
\begin{itemize}
\item[{\rm a)}] $(G,H)$ is an equal order pair.
\item[{\rm b)}] $(G,H)$ is a Camina pair.
\item[{\rm c)}] $G$ is a Frobenius group with kernel $H$.
\item[{\rm d)}] $H$ is an isolated subgroup of $G$.
\item[{\rm e)}] $d=n$ and $m_1\nmid r^{n_1}-1$ for all proper divisors $m_1$ of $m$ and $n_1$ of $n$.
\end{itemize}
\end{lemma}

\begin{proof}
\begin{description}
\item[{\rm a)} $\Leftrightarrow$ {\rm b)}] Assume that $(G,H)$ is an equal order pair and let $b^xa^y\in G$, where $1\leq x<n$ and $0\leq y<m$. Since $b^xa^y\not\in H$, we have $o(b^xa^y)=o(b^x)$, which implies that $|\langle b^xa^y\rangle|=|\langle b^x\rangle|$. Then there are $0\leq u<n$ and $0\leq v<m$ such that $\langle b^xa^y\rangle=\langle b^x\rangle^{\alpha(u,v)}$. This leads to $b^xa^y=(b^{xz})^{\alpha(u,v)}$ for some $0\leq z<o(b^x)$ with $(z,o(b^x))=1$. It follows that
\begin{equation}
b^xa^y=b^{xz}a^{v(1-r^{xz})}\nonumber    
\end{equation}and so $b^x=b^{xz}$, i.e. $z=1$. Then $b^xa^y$ is conjugate to $b^x$ and thus $(G,H)$ is a Camina pair.
\item[{\rm b)} $\Leftrightarrow$ {\rm c)}] It follows by \cite{6}, since $|H|$ and $[G:H]$ are coprime.
\item[{\rm c)} $\Leftrightarrow$ {\rm d)}] It suffices to observe that if there is $1\leq y<m$ with
\begin{equation}
a^y\in\langle b^ua^v\rangle\cap H \mbox{ for some } 1\leq u<n \mbox{ and } 0\leq v<m,\nonumber 
\end{equation}then 
\begin{equation}
b^ua^v\in C_G(a^y)\leq H,\nonumber 
\end{equation}a contradiction.
\item[{\rm d)} $\Leftrightarrow$ {\rm e)}] For every triple $(m_1,n_1,s)\in L$, we have $H_{(m_1,n_1,s)}\cap H=\langle a^{m_1}\rangle$. Then $H$ is an isolated subgroup of $G$ if and only if $H_{(m_1,n_1,s)}\cap H=1$ for all triples $(m_1,n_1,s)\in L$ satisfying $n_1\neq n$ and $\frac{m}{m_1}\mid r^{n_1}-1$, i.e. if and only if
\begin{equation}
n_1\neq n \mbox{ and } \frac{m}{m_1}\mid r^{n_1}-1 \mbox{ imply that } m_1=m.\nonumber 
\end{equation}Obviously, this is equivalent to the condition e).
\end{description}
\end{proof}


\begin{thebibliography}{99}
\bibitem[6]{6} A.R. Camina, {\it Some conditions that almost characterize Frobenius groups}, Israel J. Math. {\bf 31} (1978), 153-160.
\end{thebibliography}
\end{document}
